\documentclass[11pt]{article}
\usepackage[a4paper,margin=24mm]{geometry}
\usepackage{amsmath,amssymb,amsthm,booktabs,array}
\usepackage[T1]{fontenc}
\usepackage{lmodern,microtype}
\usepackage[colorlinks=true,linkcolor=blue,urlcolor=blue,citecolor=blue]{hyperref}
\usepackage{fancyhdr}
\pagestyle{fancy}\fancyhf{}
\fancyhead[L]{Conjecture 00000002161}
\fancyhead[R]{C0ldSmi1e}
\fancyfoot[C]{\thepage}
\setlength{\headheight}{14pt}
\setlength{\parindent}{0pt}
\setlength{\parskip}{5pt}
\newcommand{\lean}[1]{\texttt{\detokenize{#1}}}
\newcommand{\ip}{i_{+}}
\newtheorem*{result}{Counterexample}
\title{A planar tight signed graph\\\large Disproof of conjecture 00000002161}
\author{C0ldSmi1e}
\date{5 October 2026}
\begin{document}
\maketitle
\begin{abstract}
The all-positive four-vertex path has independence number $2$, positive
inertia index $2$, and minimum orientable genus $0$. It is connected and
noncomplete. It therefore contradicts the asserted minimum genus $1$,
whether the minimum is taken separately for every eligible graph or across
the eligible class. A Lean 4 project verifies the counterexample using the
standard rotation-system definition of orientable cellular embeddings.
\end{abstract}

\section{Statement and scope}
The exact bilingual source is supplied as \lean{conjecture.md}. Its English
text is:
\begin{quote}\small
Definition: The inertia bound of signed graphs states the independence number
is at most the positive inertia index. Conjecture: The minimal genus of graphs
attaining the tight bound equals 1 (apart from complete graphs); the genus is
computed explicitly via the Euler formula of 2-cell embeddings.
\end{quote}
We retain the printed equality condition $\alpha(G)=\ip(A)$, where $A$ is the
signed adjacency matrix. We neither replace the printed inequality by a
conventional inertia bound nor assume it holds for every signed graph.
The witness satisfies the printed equality itself. Genus has its usual
orientable meaning; a two-cell embedding has open disks as complementary
faces. This submission makes no claim about a reformulation requiring a
strictly nonorientable surface.

\section{A tight, connected, noncomplete witness}
Take the path $0-1-2-3$, with all three edge signs $+1$. The vertices $0,2$
are nonadjacent, so the graph is not complete. The path is connected.
The set $\{0,2\}$ is independent, while the disjoint edges $01$ and $23$
cover all vertices and each contributes at most one member to an independent
set. Thus $\alpha(P_4)=2$.

The adjacency matrix and its characteristic polynomial are
\[
 A=\begin{pmatrix}0&1&0&0\\1&0&1&0\\0&1&0&1\\0&0&1&0\end{pmatrix},
 \qquad p(t)=\det(tI-A)=t^4-3t^2+1.
\]
Put $r=(\sqrt5+1)/2$ and $s=(\sqrt5-1)/2$. Both are positive,
$rs=1$, and $r^2+s^2=3$. Hence
\[
 p(t)=(t-r)(t-s)(t+r)(t+s).
\]
Its complete root multiset is $[r,s,-r,-s]$. Counting positive eigenvalues
with algebraic multiplicity gives $\ip(A)=2$; no numerical approximation is
used. This is the ordinary inertia convention for real symmetric matrices
\cite{inertia}.

\newpage
\section{An actual cellular embedding and its genus}
\subsection{Rotation systems and their meaning}
For a connected finite graph with no isolated vertex, use one dart for each
ordered adjacent pair. Let $\lambda$ reverse darts and let $\rho$ cyclically
order the darts with each fixed tail. The formal data require that reversal
is a fixed-point-free involution, that the dart-to-ordered-edge map is
bijective, and that the rotation preserves tails and has one orbit in each
vertex star. These conditions describe the original graph, not a different
multigraph or an arbitrary pair of permutations.

This is the standard combinatorial encoding of an orientable cellular
embedding \cite[Section 2, Proposition 2.2]{rotation}. Its geometric meaning
can be seen by taking oriented vertex disks, joining their marked boundary
intervals with untwisted edge bands, and capping the resulting boundary
circles with disks. The boundary components correspond to cycles of
$\rho\lambda$, with $\lambda$ applied first. Conversely, an oriented
cellular embedding gives the local cyclic orders.

For this path specifically, begin with the vertex-$0$ disk and attach each
successive band and new vertex disk. Each step joins two disks along a band
and produces a disk. The thickening is therefore a disk; capping its one
boundary circle gives the sphere. The central edge arcs have the prescribed
endpoints and disjoint interiors. Thus a genuine genus-zero cellular
embedding exists, independently of merely guessing a genus from numbers.

\subsection{The complete finite certificate}
The following table fixes all dart data; entries in the last two rows are
dart labels.
\begin{center}
\begin{tabular}{c*{6}{c}}
\toprule
Dart $d$ & 0 & 1 & 2 & 3 & 4 & 5\\
\midrule
Directed edge & $01$ & $10$ & $12$ & $21$ & $23$ & $32$\\
$\lambda(d)$ & 1 & 0 & 3 & 2 & 5 & 4\\
$\rho(d)$ & 0 & 2 & 1 & 4 & 3 & 5\\
\bottomrule
\end{tabular}
\end{center}
Thus $\lambda=(0\ 1)(2\ 3)(4\ 5)$,
$\rho=(0)(1\ 2)(3\ 4)(5)$, and
\[
 \rho\lambda=(0\ 2\ 4\ 5\ 3\ 1).
\]
There is one face, with boundary walk $0,1,2,3,2,1,0$. Repeated edges in this
boundary are allowed: each bridge borders the same face on both sides.
The vertex, edge, and face counts are $V=4$, $E=3$, $F=1$. Euler's formula
for this orientable cellular embedding gives
\[
 2-2g=V-E+F=4-3+1=2,\qquad g=0.
\]
The minimum genus is also $0$: it is attained here and every orientable
genus is a nonnegative integer.

\begin{result}
The signed path $P_4$ is connected and noncomplete, attains
$\alpha(P_4)=\ip(A)=2$, and has minimum genus $0$.
\end{result}
If the conjecture asserts minimum genus $1$ for every tight noncomplete
graph, this witness contradicts it directly. If it asserts a class minimum
of $1$, it necessarily asserts a lower bound $1$ for each eligible graph's
minimum genus; the same witness contradicts that clause. Restricting the
formal universal clauses to connected graphs is sufficient, since a
counterexample in that subclass is a counterexample in the full class.

\newpage
\section{Formal specification and coverage}
The project uses Lean 4.19.0 and Mathlib v4.19.0, with the exact dependency
revisions recorded in \lean{lean/lake-manifest.json}. All statements below
are in namespace \lean{Conjecture2161}.

\begin{center}\small
\begin{tabular}{@{}>{\raggedright\arraybackslash}p{.42\textwidth}>{\raggedright\arraybackslash}p{.53\textwidth}@{}}
\toprule
Definition or theorem & What is represented or proved\\
\midrule
\lean{SignedGraph} & An actual simple graph, symmetric real matrix, edge
entries $\pm1$, and zero entries on every nonedge.\\[3pt]
\lean{independenceNumber} & The maximum cardinality over all finite independent
subsets, with its generic upper and lower specifications proved.\\[3pt]
\lean{path_independence_number} & The witness set and upper bound over all
16 subsets establish the exact value $2$.\\[3pt]
\lean{path_roots},\newline \lean{path_positive_inertia} & The actual matrix
characteristic polynomial has the complete root multiset above; its positive
part has cardinality $2$.\\[3pt]
\lean{RotationSystem} & Every incidence, reversal, graph-connectedness,
nonempty-star and cyclic-order condition described on page 2.\\[3pt]
\lean{RotationSystem.Face} & The quotient of darts by the actual face-permutation
orbit relation. A named instance proves this type finite.\\[3pt]
\lean{HasCellularGenus} & Existence of a valid rotation system satisfying
$V-E+F=2-2g$ with integer arithmetic.\\[3pt]
\lean{IsMinimumGenus} & A realized genus and a lower bound for every other
realized genus; no empty infimum or default genus value.\\[3pt]
\lean{counterexample_certificate} & Connectedness, noncompleteness, exact
printed equality, and minimum genus $0$.\\[3pt]
\lean{not_per_graph_genus_one},\newline \lean{not_class_minimum_one} & Negations
of the necessary clauses for the two readings of the conjecture.\\
\bottomrule
\end{tabular}
\end{center}

\paragraph{Formalization boundary.}
Lean checks the entire finite rotation-system certificate, actual face
quotient cardinality, integer Euler equality, existential genus certificate,
minimum argument, and final negations. The project uses the standard
combinatorial definition of orientable cellular embedding, with the
geometric correspondence explained above and cited. It does not additionally
formalize the disk-and-band realization theorem in Mathlib's general
topology language. No geometric-realization axiom is inserted into Lean.
The restriction to connected noncomplete graphs avoids the isolated-vertex
degeneracy of a dart-only representation.

\paragraph{Trust and execution.}
The 50 handwritten declarations comprise 20 definitions/structures,
27 theorems, and 3 named instances. \lean{Audit2161.lean} prints the axiom
dependencies of all 50; \lean{Check.lean} also displays the definitions and
theorem types. The project contains no admitted proof, \lean{native_decide},
or custom axiom. Small finite checks use ordinary kernel-checked
\lean{decide}. Builds and source replays treat warnings as errors.
An additional inspection inventories generated declarations and checks
transitive logical dependencies. Every logical declaration depends only on
\lean{propext}, \lean{Classical.choice}, and \lean{Quot.sound}. Compiler
runtime artifacts are inventoried separately and are absent from all
logical dependency closures.

\paragraph{The printed inequality.}
The purported universal inequality $\alpha\leq\ip$ is not used as an axiom:
for example, $P_3$ has $\alpha=2$ and $\ip=1$. That observation is only a
scope check. The disproof above rests on $P_4$, which attains the equality
exactly as printed, together with its genus-zero embedding.

\newpage
\section{Independent computations and reproduction}
\subsection{Exact auxiliary checks}
The separate Python checker in \lean{auxiliary/} uses exact rational
arithmetic and only the standard library. It enumerates all 16 vertex
subsets and expands the determinant by all 24 permutations. For $p(t)=t^4-3t^2+1$, its exact Sturm sequence is
\[
 p(t),\quad 4t^3-6t,\quad \tfrac32t^2-1,\quad \tfrac{10}{3}t,\quad 1.
\]
The sign-variation counts at $-\infty,0,+\infty$ are $4,2,0$, giving two
negative and two positive roots and no zero root. The checker handles
multiplicity through square-free layers; a repeated-root test is included.
The Sturm convention is documented in \cite{sturm}.

The same checker independently verifies every dart, vertex rotation and
face orbit. It realizes the face as an oriented six-sided polygon with
paired sides, reconstructs the quotient vertices, and verifies each vertex
link is a circle. It also checks the straight-line drawing with vertices
$(0,0),(1,0),(2,0),(3,0)$ for edge crossings and nonincident vertices in
edges. Boundary tests include $K_2$, $P_3$, $K_3$, and deliberately invalid
certificates/drawings. The supplied certificate records all intermediate
values, including $(V,E,F)=(4,3,1)$ and $g=0$.

\subsection{Commands}
From the submission folder, with Python 3 and the pinned Lean toolchain:
\begin{verbatim}
cd lean
lake exe cache get
lake build
lake env lean -DwarningAsError=true Conjecture2161.lean
lake env lean -DwarningAsError=true Audit2161.lean
lake env lean -DwarningAsError=true Check.lean
cd ..
python3 auxiliary/replay.py
\end{verbatim}
The pinned manifest must be preserved. Initial dependency/cache retrieval
requires network access. The replay runs the exact checker afresh and
requires byte-for-byte agreement with the supplied certificate.
\lean{README.md} and \lean{verification/} document the independent clean-build
procedure, dependency identities, source hashes, trust audit, and local
semantic review. These are local verification records, not maintainer
acceptance or a replacement for the repository's review process.

\begin{thebibliography}{9}\small\raggedright
\bibitem{rotation} J. L. Gross and T. W. Tucker,
\emph{Enumerating Graph Embeddings and Partial-Duals by Genus and Euler Genus},
Enumerative Combinatorics and Applications 1:1 (2021), Article S2S1.
\href{https://doi.org/10.54550/ECA2021V1S1S1}{doi:10.54550/ECA2021V1S1S1}.
\bibitem{inertia} Y. Nakatsukasa and V. Noferini,
\emph{Inertia laws and localization of real eigenvalues for generalized
indefinite eigenvalue problems}.
\href{https://arxiv.org/abs/1711.00495}{arXiv:1711.00495}.
\bibitem{sturm} P. P\'ebay, J. M. Rojas, and D. C. Thompson,
\emph{Sturm's Theorem with Endpoints}.
\href{https://arxiv.org/abs/2208.07904}{arXiv:2208.07904}.
\end{thebibliography}
\end{document}
